\section{Conclusion}
\label{sec:conclusion}

The main contribution of this thesis is a set of fully operational forecasting
models for load, solar, and wind generation based exclusively on open data, the
OBTF models. They are released as part of an open pipeline that enforces
forecast-time information availability and that others can rerun and reuse.
Their daily submission to the Energy-Arena benchmark
\cite{kleinebrahm2026energyarena} lets anyone verify their accuracy beyond
the test period of this thesis. In
contrast to the \mbox{ENTSO-E} forecasts, whose methods are not disclosed, their
inputs and construction are fully open. Each model is either more accurate than its
\mbox{ENTSO-E} counterpart, as for load, or admissible before gate closure when
the \mbox{ENTSO-E} forecast is not, as for solar and wind. Used as explicit
inputs, the OBTF forecasts improve the operational day-ahead price forecast for
the \mbox{DE-LU} bidding zone, in which every input must be available before the
12:00 gate closure. The models can therefore be reused as admissible,
transparent inputs in other operational applications.

For RQ1, the OBTF load forecast is more accurate than the published
\mbox{ENTSO-E} forecast in the operational setting, reducing the load RMSE by
36.0\%. The OBTF solar and wind forecasts are less accurate than the
\mbox{ENTSO-E} forecasts, but those benchmarks are not published before gate
closure and cannot be used in an operational price forecast.

For RQ2, adding the OBTF load, combined renewable-generation, and
residual-load forecasts to the price model lowers the day-ahead price RMSE by
11.2\%, a statistically significant improvement. The weather-free variant is
statistically indistinguishable in accuracy, so the raw weather fields add no detectable accuracy once these
forecasts are included, which makes them a compact representation of the
weather-driven system state for the price model.

For RQ3, price accuracy improves as the morning progresses, but gradually. Up to 11:00 the error falls by 5.6\% without a statistically significant step,
which may partly reflect that the backtest uses the same weather run at all
cutoffs. The only significant gain follows at gate closure once the EXAA price
becomes admissible. That
the earlier-clearing EXAA price captures most of the achievable accuracy at 12:00
is a known property of this market rather than a result of this thesis
\cite{eiser2026operational}. Before any market price is available, the forecast
rests on price lags, weather, the OBTF forecasts, and the
reserve-market results.

The findings are relevant beyond the market and period studied here. Day-ahead
prices depend strongly on weather-driven load and renewable generation
\cite{sensfuss2008meritOrder,pape2016fundamentals}, and the corresponding public
forecasts are either less accurate than the OBTF forecast or unavailable before
gate closure. Turning open weather data into admissible load and
renewable-generation forecasts is therefore valuable in other bidding zones with
the same gate-closure constraint. In a field where forecasts are often produced by
undisclosed models and evaluated retrospectively, such an open and continuously
benchmarked pipeline offers a transparent operational baseline that others can
reproduce, extend, and compare against. The thesis is thus a step toward
day-ahead price forecasting that is operational, open, and explicit about what
information it uses and when.
